% ============================================================
\chapter{Conclusiones y Trabajo Futuro}
\label{cap:conclusiones}
% ============================================================

\section{Conclusiones}

Este trabajo ha demostrado la viabilidad de un enfoque híbrido de tres capas para la
detección automática de cadenas de ataque en infraestructuras \textit{Kubernetes}. Las
conclusiones C1--C6 responden, respectivamente, a los objetivos OE2, OE4, OE1, OE5,
OE3 y OE6 definidos en la Sección~\ref{cap:objetivos}. Las conclusiones principales son:

\textbf{C1 — La Capa~1 supera ampliamente el objetivo de diseño (OE2).} El clasificador
\textit{Random Forest} alcanza F1 macro\,=\,0,9935 y \ac{AUC}\,=\,0,9988 en test, superando
el umbral objetivo del 0,85 en 14,4 puntos porcentuales. El modelo de severidad de
tres clases logra F1\,=\,1,000 en la clase \textit{high-critical}, garantizando que
ninguna misconfiguration crítica es perdida.

\textbf{C2 — La Capa~2 supera el objetivo de Precisión@5 (OE4).} La \ac{GAT} de
3~capas con \textit{embedding} de tipos de arista alcanza
F1~macro\,=\,0,870\,\(\pm\)\,0,075 y P@5\,=\,0,720\,\(\pm\)\,0,098 en los 5~pliegues
bajo el protocolo \textit{group-aware}, superando el objetivo de P@5\,\(>\)\,0,70.
Tres pliegues alcanzan su techo teórico de P@5, situando todos los clústeres de
cadena del pliegue en las primeras 5~posiciones del ranking. La ablación de
arquitectura (Tabla~\ref{tab:arch_ablation}) respalda las decisiones de diseño: la
\ac{GAT} supera a la \ac{GCN} en 8~puntos de P@5, y la configuración de 3~capas es
la única que cumple el objetivo de la métrica primaria.

\textbf{C3 — La combinación de aumentación de datos y expansión del corpus es la
contribución metodológica central (OE1).} La evolución en cinco fases — línea base
(P@5\,=\,0,40), aumentación (P@5\,=\,0,52), corrección \texttt{HOSTPATH\_MOUNT}
(P@5\,=\,0,60), dataset extendido con 25~cadenas (P@5\,=\,0,88, protocolo
preliminar) y corrección \textit{group-aware} del particionado (P@5\,=\,0,72,
resultado final) — demuestra que la escasez de ejemplos de cadena de ataque es el
factor limitante predominante y que ambas estrategias de mitigación son efectivas y
complementarias. La quinta fase ilustra además una lección metodológica: la
aumentación de datos exige particionado consciente de grupos, pues parte de la
mejora aparente de la fase~4 procedía de fuga entre variantes aumentadas y sus
grafos base de evaluación.

\textbf{C4 — El ensemble es robusto y no produce falsos positivos de clústeres
limpios (OE5).} La ablación sobre el conjunto de test
(Tabla~\ref{tab:ensemble_ablation}) confirma que el componente \ac{GNN} es
indispensable: la configuración RF-sólo degrada el ranking (P@5 de 0,60 a 0,40) e
introduce un falso positivo de clúster limpio (\ac{FPR}\,=\,0,20), mientras que
\ac{FPR}\_clean\,=\,0,00 para cualquier combinación que incluya el \ac{GNN}. El
clúster de mayor puntuación es una cadena de ataque (P@1\,=\,1,00). Con honestidad
metodológica debe señalarse que, en este conjunto de test, los pesos optimizados
por el \ac{GA} no superan a los pesos uniformes ni al \ac{GNN} en solitario: la
aportación del \textit{ensemble} en estos datos es la robustez frente a la
degradación de componentes individuales, no una mejora del ranking sobre el mejor
componente.

\textbf{C5 — El tipo de arista \ac{RBAC} mejora la cobertura de escalada de privilegios
(OE3).} La detección de roles privilegiados en dos pasadas sobre los manifiestos,
limitando las aristas \ac{RBAC} a cuentas de servicio vinculadas realmente a roles con
permisos de escalada, produce señales de grafo más precisas y reduce los falsos
positivos estructurales.

\textbf{C6 — kubescan es una herramienta distribuible, instalable y operacionalmente
viable (OE6).} La evaluación funcional de la Sección~\ref{sec:evaluacion_funcional}
confirma que el sistema se instala sin compiladores externos (\textit{pip install -e
kubescan/}), analiza clústeres de hasta 312~manifiestos en menos de 8~segundos en
hardware de gama media, produce salida en dos formatos (texto y \ac{JSON}) adecuados para
integración en pipelines CI/CD, y garantiza reproducibilidad completa mediante la
resolución determinista de \textit{checkpoints}. El prototipo implementado en este
trabajo demuestra que el paradigma \textit{shift-left} —detectar cadenas de ataque
antes del despliegue, sin requerir acceso al clúster activo— es técnicamente factible
y operacionalmente eficiente.

\section{Consideraciones éticas y uso responsable}
\label{sec:etica}

\textit{kubescan} es una herramienta de uso dual: aunque su propósito es defensivo
—detectar cadenas de ataque potenciales antes de que sean explotadas— la información
que produce (identificación de pods con flags de escape, rutas de privilegio \ac{RBAC})
podría ser utilizada con fines ofensivos si cayera en manos de un actor no autorizado.
En consecuencia, este trabajo adopta los principios de divulgación responsable
(\textit{responsible disclosure}) que rigen la investigación en ciberseguridad.

Los patrones de ataque detectados por \textit{kubescan} se basan exclusivamente en
documentación pública —la guía \ac{NSA}/\ac{CISA} \citep{NSA2022}, el trabajo de BishopFox
\citep{BishopFox2021} y el marco MITRE ATT\&CK for Containers \citep{MITRE2021Containers}—
y no revelan ninguna vulnerabilidad nueva. El código fuente del sistema es de acceso
abierto para facilitar la auditoría independiente de las reglas de detección, lo que
permite a la comunidad verificar que las señales de riesgo generadas son correctas y
no contienen sesgos perjudiciales. El uso de la herramienta sobre clústeres de terceros
sin autorización explícita estaría fuera del ámbito previsto y podría constituir un
acceso no autorizado a sistemas informáticos según la legislación aplicable.

\section{Limitaciones}
\label{sec:limitaciones}

\textbf{L1 — Varianza entre pliegues.} Con 81~grafos originales en los pliegues de
validación cruzada, la varianza entre pliegues (\(\sigma\)\,=\,0,075 para F1 macro;
rango 0,738--0,941) refleja la dependencia de los resultados respecto a qué
clústeres caen en cada pliegue de validación. Un corpus más amplio reduciría esta
varianza.

\textbf{L2 — Imputación de características extendidas.} Las 6 características
extendidas (basadas en \textit{Checkov}) están imputadas para el 42,6\,\% de las
filas, diluyendo su señal. La descarga completa de los manifiestos referenciados en
el dataset de Rahman et~al.\ permitiría mejorar estos vectores. Un experimento de
ablación eliminando las 6 características imputadas y reentrenando únicamente con las
19 características de alta cobertura podría cuantificar el impacto de esta limitación
sobre la F1 macro de la Capa~1.

\textbf{L3 — Sesgo de muestreo en el dataset tabular.} Los 10 repositorios más
representados suponen el 60,8\,\% del dataset, y \textit{gitcontroller} por sí solo
aporta el 13,7\,\%. Un esquema de validación cruzada por repositorio ofrecería
estimaciones más conservadoras de la generalización del modelo.

\textbf{L4 — Evaluación en clústeres de producción reales.} Toda la evaluación
se ha realizado sobre repositorios públicos. La validación en entornos de producción
de organizaciones reales es la siguiente etapa crítica para confirmar la validez
externa del sistema.

\textbf{L5 — Identificabilidad de los pesos del ensemble.} La Capa~3 optimiza
los pesos \((w_{\text{rf}}, w_{\text{gnn}}, w_{\text{esc}})\) maximizando la función
objetivo \(F = 0{,}7 \cdot P@5 + 0{,}3 \cdot (1 - \text{\ac{FPR}\_clean})\) sobre las
predicciones \textit{out-of-fold} (\acs{OOF}) de los 5~modelos de validación
cruzada, calculadas sobre los 81~grafos originales de entrenamiento/validación
(21~cadenas); bajo el protocolo \textit{group-aware}, ningún clúster de test
participa en este ajuste. La limitación restante es de identificabilidad: con una
función objetivo basada en P@5 —que solo toma seis valores discretos— y 81~muestras,
el paisaje de \textit{fitness} presenta amplias mesetas en las que múltiples
combinaciones de pesos alcanzan un objetivo idéntico. La búsqueda en rejilla y el
\ac{GA} encuentran configuraciones muy distintas con la misma puntuación, y el
\ac{GA} converge hacia pesos casi uniformes
\((w_{\text{rf}} \approx 0{,}34,\, w_{\text{gnn}} \approx 0{,}33,\,
w_{\text{esc}} \approx 0{,}33)\) — un punto compatible con la deriva hacia el
centroide que induce el cruce aritmético en mesetas planas. Los pesos concretos
deben interpretarse, por tanto, como una solución representativa de una clase de
equivalencia, no como un óptimo único; el tamaño del conjunto de test (15~grafos)
limita además la significación estadística de la comparación entre configuraciones
de pesos, como reflejan los intervalos de confianza reportados.

\textbf{L6 — Etiquetado derivado de reglas sobre el mismo espacio de
características.} Las etiquetas de grafo (clean / isolated / attack\_chain) se
generan mediante una regla determinista definida sobre las mismas \textit{flags} y
aristas que observan los modelos (Sección de diseño). El sistema aprende, por
tanto, una aproximación robusta de dicha regla bajo ruido, enmascaramiento y
observabilidad parcial — no una noción de cadena de ataque independiente de la
regla. Esta circularidad es común en la literatura de detección de
misconfiguraciones, pero limita el alcance de la afirmación «predice cadenas de
ataque». Las vías para romperla son la validación externa frente a herramientas de
grafo de ataque sobre clústeres activos (KubeHound, IceKube), el etiquetado experto
de una submuestra, y la confirmación mediante explotación controlada en entornos de
laboratorio como \textit{kubernetes-goat}.

\section{Trabajo futuro}

Los resultados de este trabajo permiten formular las siguientes recomendaciones de
trabajo futuro, ordenadas por impacto esperado:

\textbf{TF1 — Continuar la expansión del corpus.} El objetivo de P@5\,\(>\)\,0,70 ha
sido alcanzado con 25~clústeres de cadena de ataque. La incorporación de corpus
adicionales — repositorios de infraestructura de organizaciones de código abierto,
entornos de laboratorio multiclúster — permitiría reducir la varianza entre pliegues
y mejorar la generalización a patrones de ataque menos representados.

\textbf{TF2 — Implementar P@5 global como métrica alternativa.} En lugar de P@5 por
pliegue, evaluar el ranking del modelo sobre el conjunto completo de grafos en un único
ranking global eliminaría el efecto de techo por pliegue y permitiría comparar con el
objetivo original de diseño sin restricciones estructurales.

\textbf{TF3 — Incorporar análisis semántico de \ac{RBAC} completo.} El parser de \ac{YAML}
actual identifica roles privilegiados por nombre y por verbos de escalada. Una extensión
natural es el análisis de la cadena completa Role/ClusterRole \(\to\) RoleBinding \(\to\)
ServiceAccount \(\to\) Pod para detectar rutas de escalada de privilegios mediante
\ac{RBAC} que no requieran flags de escape directas.

\textbf{TF4 — Validación cruzada dejando un repositorio fuera.} La estrategia
\textit{leave-one-repo-out} ofrece estimaciones más conservadoras de la generalización
que la validación cruzada por fila, y es especialmente relevante dada la concentración
del dataset en pocos repositorios.

\textbf{TF5 — Modo de análisis continuo (\textit{admission controller}).} La integración
de \textit{kubescan} como un \textit{validating admission webhook} de \textit{Kubernetes}
permitiría bloquear el despliegue de manifiestos que eleven el riesgo de cadena de
ataque del clúster en tiempo real, sin necesidad de análisis \textit{post-hoc}.

\textbf{TF6 — Estudio de ablación arquitectónica.} La justificación de la elección
de \ac{GAT} frente a \ac{GCN}, de 3~capas frente a 2, y de \textit{mean+max pooling} frente a
\textit{mean-only} es actualmente teórica (fundamentada en la literatura de \ac{GNN}s
para grafos heterogéneos, Sección~\ref{sec:gnn_por_clase}). Un experimento de ablación
controlado que varíe cada uno de estos hiperparámetros arquitectónicos manteniendo
constante el resto cuantificaría la contribución incremental de cada decisión de diseño
y produciría evidencia empírica directamente comparable con la literatura de \ac{GNN}s para
ciberseguridad.
